\section{Conclusion}
\label{sec:latent_space_conclusion}

% figures/latent_space/latent_space_overview.tex

\providecommand{\lsePanelWidth}{5.5cm}
\newlength{\ovPanelWd}
\newlength{\ovPanelHt}
\newlength{\ovPanelGap}
\newlength{\ovRowGap}

\setlength{\ovPanelWd}{\lsePanelWidth}
\pgfmathsetlength{\ovPanelHt} {\ovPanelWd * 5 / 6}
\pgfmathsetlength{\ovPanelGap}{\ovPanelWd * 1}
\pgfmathsetlength{\ovRowGap}  {\ovPanelHt * 1}

\begin{figure}[t]
\centering

% ================================================================
%  latent_space_shared.tex
%  Couleurs, macros et styles partagés par toutes les figures
%  latent_space_*.tex
%  Usage : % ================================================================
%  latent_space_shared.tex
%  Couleurs, macros et styles partagés par toutes les figures
%  latent_space_*.tex
%  Usage : \input{figures/latent_space/latent_space_shared.tex}
% ================================================================

% ----------------------------------------------------------------
% Couleurs des clusters
% ----------------------------------------------------------------
\colorlet{colorA}{blue!70!cyan}
\colorlet{colorB}{orange}
\colorlet{colorC}{green!60!black}

% Couleurs de directions (latent_space_directions.tex)
\colorlet{colorDirBA}{violet}
\colorlet{colorDirBC}{red!70!black}

% ----------------------------------------------------------------
% Macros de dessin
% ----------------------------------------------------------------
\input{figures/latent_space/plot_points.tex}
\input{figures/latent_space/plot_ellipse.tex}
\input{figures/latent_space/plot_probe.tex}
\input{figures/latent_space/plot_direction.tex}

% ----------------------------------------------------------------
% Styles d'axe partagés
% ----------------------------------------------------------------
\pgfplotsset{
  % Axe standard (clustering, probing, directions)
  latent axis/.style={
    width=11cm,
    height=10cm,
    xmin=-7, xmax=8,
    ymin=-7, ymax=5.5,
    axis x line=bottom,
    axis y line=left,
    axis line style={
      -{Stealth[length=8pt, width=6pt]},
      thick
    },
    ticks=none,
    xlabel={},
    ylabel={},
    clip=true,
  },
  % Variante panneau carré (latent_space_sparse_autoencoders.tex)
  latent panel/.style={
    latent axis,
    width=9cm,
    height=9cm,
  },
  % Variante petit panneau (latent_space_embeddings.tex)
  % N.B. width/height sont passés explicitement par chaque fichier
  % via width=\lsePanelWd, height=\lsePanelHt dans l'axe.
  latent embed/.style={
    latent axis,
    every axis plot/.append style={forget plot},
  },
}

% ================================================================

% ----------------------------------------------------------------
% Couleurs des clusters
% ----------------------------------------------------------------
\colorlet{colorA}{blue!70!cyan}
\colorlet{colorB}{orange}
\colorlet{colorC}{green!60!black}

% Couleurs de directions (latent_space_directions.tex)
\colorlet{colorDirBA}{violet}
\colorlet{colorDirBC}{red!70!black}

% ----------------------------------------------------------------
% Macros de dessin
% ----------------------------------------------------------------
% ================================================================
%  \plotPoints — Reusable macro to render a point cloud from file
%
%  Usage:
%    \plotPoints[color=<c>, size=<s>, opacity=<o>]{<file.data>}
%
%  Arguments (all optional, defaults shown below):
%    color   : fill & stroke color of the markers  (default: black)
%    size    : marker radius                        (default: 2pt)
%    opacity : fill opacity [0..1]                  (default: 0.85)
%
%  Positional (mandatory):
%    #1 : path to a whitespace-separated x y data file
%         (lines starting with % are treated as comments)
%
%  Example:
%    \plotPoints[color=blue!70!cyan, size=3pt]{data/cluster_a.data}
% ================================================================
\pgfkeys{
  /plotPoints/.is family, /plotPoints,
  % --- defaults ---
  color/.initial   = black,
  size/.initial    = 2pt,
  opacity/.initial = 0.85,
}

\newcommand{\plotPoints}[2][]{%
  \pgfkeys{/plotPoints, #1}%                     % apply user overrides
  \pgfkeysgetvalue{/plotPoints/color}  {\ppColor}%
  \pgfkeysgetvalue{/plotPoints/size}   {\ppSize}%
  \pgfkeysgetvalue{/plotPoints/opacity}{\ppOpacity}%
  \addplot[
    only marks,
    mark        = *,
    mark size   = \ppSize,
    mark options= {fill=\ppColor, draw=\ppColor, fill opacity=\ppOpacity},
  ] file {#2};%
}
% ================================================================
%  \plotEllipse — Companion macro to draw a confidence ellipse
%
%  Usage:
%    \plotEllipse[color=<c>, width=<w>]
%                {<cx>}{<cy>}{<angle>}{<rx>}{<ry>}
%
%  Arguments (all optional):
%    color : stroke color  (default: black)
%    width : line width    (default: very thick)
%
%  Positional (mandatory):
%    {cx}    : ellipse center x  (axis coordinates)
%    {cy}    : ellipse center y  (axis coordinates)
%    {angle} : counter-clockwise rotation in degrees
%    {rx}    : horizontal semi-axis (cm)
%    {ry}    : vertical   semi-axis (cm)
%
%  Example:
%    \plotEllipse[color=orange]{-3}{1.5}{15}{2.3}{1.2}
% ================================================================
\pgfkeys{
  /plotEllipse/.is family, /plotEllipse,
  color/.initial = black,
  width/.initial = very thick,
}

\newcommand{\plotEllipse}[6][]{%
  \pgfkeys{/plotEllipse, #1}%
  \pgfkeysgetvalue{/plotEllipse/color}{\peColor}%
  \pgfkeysgetvalue{/plotEllipse/width}{\peWidth}%
  \draw[\peColor, \peWidth, rotate around={#4:(#2,#3)}]
    (#2,#3) ellipse (#5cm and #6cm);%
}
% ================================================================
%  plot_probe.tex — Linear probing decision boundary
%
%  \plotProbe[<options>]{<slope>}{<intercept>}
%
%  Draws  y = slope*x + intercept  clipped to the axis bounds.
%
%  Options (all optional):
%    color : line color   (default: black)
%    style : dash pattern (default: dashed)
%    width : line width   (default: thick)
%
%  WHY \edef ON THE WHOLE CALL: same deferred-rendering issue as
%  plot_points.tex — pgfplots evaluates ALL \addplot options at
%  \end{axis}, not at call time, so macro values must be baked in.
%
%  Example:
%    \plotProbe[color=colorA]{1}{3}
% ================================================================
\pgfkeys{
  /plotprobe/.is family, /plotprobe,
  color/.initial = black,
  style/.initial = dashed,
  width/.initial = thick,
}

\newcommand{\plotProbe}[3][]{%
  \pgfkeys{/plotprobe, #1}%
  \pgfkeysgetvalue{/plotprobe/color}{\prColor}%
  \pgfkeysgetvalue{/plotprobe/style}{\prStyle}%
  \pgfkeysgetvalue{/plotprobe/width}{\prWidth}%
  \edef\prDoPlot{%
    \noexpand\addplot[%
      domain=-7:8, samples=2,%
      \prStyle, \prWidth, color=\prColor, forget plot%
    ] {#2 * x + (#3)};%
  }%
  \prDoPlot
}

% ================================================================
%  plot_direction.tex — Interpretable direction vector
%
%  \plotDirection[<options>]{<x1>}{<y1>}{<x2>}{<y2>}
%
%  Draws a thick arrow from (x1,y1) to (x2,y2) in axis coordinates.
%  Must be called OUTSIDE the \addplot system, directly inside axis.
%
%  Options (all optional):
%    color   : arrow color   (default: black)
%    width   : line width    (default: very thick)
%    shorten : shorten tip   (default: 4pt)
%
%  Example:
%    \plotDirection[color=colorA]{-2.5}{-4}{-3}{1.5}
% ================================================================
\pgfkeys{
  /plotdirection/.is family, /plotdirection,
  color/.initial   = black,
  width/.initial   = very thick,
  shorten/.initial = 4pt,
}

\newcommand{\plotDirection}[6][]{%
  \pgfkeys{/plotdirection, #1}%
  \pgfkeysgetvalue{/plotdirection/color}  {\pdColor}%
  \pgfkeysgetvalue{/plotdirection/width}  {\pdWidth}%
  \pgfkeysgetvalue{/plotdirection/shorten}{\pdShorten}%
  \draw[\pdColor, \pdWidth, -Stealth,
        shorten >=\pdShorten, shorten <=\pdShorten]
    (axis cs:#2,#3) -- (axis cs:#4,#5);%
}


% ----------------------------------------------------------------
% Styles d'axe partagés
% ----------------------------------------------------------------
\pgfplotsset{
  % Axe standard (clustering, probing, directions)
  latent axis/.style={
    width=11cm,
    height=10cm,
    xmin=-7, xmax=8,
    ymin=-7, ymax=5.5,
    axis x line=bottom,
    axis y line=left,
    axis line style={
      -{Stealth[length=8pt, width=6pt]},
      thick
    },
    ticks=none,
    xlabel={},
    ylabel={},
    clip=true,
  },
  % Variante panneau carré (latent_space_sparse_autoencoders.tex)
  latent panel/.style={
    latent axis,
    width=9cm,
    height=9cm,
  },
  % Variante petit panneau (latent_space_embeddings.tex)
  % N.B. width/height sont passés explicitement par chaque fichier
  % via width=\lsePanelWd, height=\lsePanelHt dans l'axe.
  latent embed/.style={
    latent axis,
    every axis plot/.append style={forget plot},
  },
}


\resizebox{\textwidth}{!}{
\begin{tikzpicture}

  % ============================================================
  % (a) TOP-LEFT — Probing
  % ============================================================
  \begin{scope}[xshift=0pt, yshift=0pt]
  \begin{axis}[
    name=axisProbing,
    latent embed, width=\ovPanelWd, height=\ovPanelHt,
  ]
    \fill[colorA, opacity=0.10]
      (axis cs:-7,-4) -- (axis cs:-7,5.5) -- (axis cs:2.5,5.5) -- cycle;
    \fill[colorB, opacity=0.10]
      (axis cs:-7,0.1) -- (axis cs:8,-4.4) --
      (axis cs:8,-7)   -- (axis cs:-7,-7)  -- cycle;
    \fill[colorC, opacity=0.10]
      (axis cs:-0.75,5.5) -- (axis cs:8,5.5) --
      (axis cs:8,-7)      -- (axis cs:5.5,-7) -- cycle;
    \plotProbe[color=colorA]{1}{3}
    \plotProbe[color=colorB]{-0.3}{-2}
    \plotProbe[color=colorC]{-2}{4}
    \plotPoints[color=colorA, mark=*]{figures/latent_space/cluster_a.data}
    \plotPoints[color=colorB, mark=square*]{figures/latent_space/cluster_b.data}
    \plotPoints[color=colorC, size=2pt, mark=triangle*]{figures/latent_space/cluster_c.data}
  \end{axis}
  \node[anchor=north, yshift=-6pt] at (axisProbing.south)
    {\textbf{(a)} Probing};
  \end{scope}

  % ============================================================
  % (b) TOP-RIGHT — Interpretable directions
  % ============================================================
  \begin{scope}[xshift=\ovPanelGap, yshift=0pt]
  \begin{axis}[
    name=axisDirections,
    latent embed, width=\ovPanelWd, height=\ovPanelHt,
  ]
    \plotPoints[color=colorA, mark=*]{figures/latent_space/cluster_a.data}
    \plotPoints[color=colorB, mark=square*]{figures/latent_space/cluster_b.data}
    \plotPoints[color=colorC, size=2pt, mark=triangle*]{figures/latent_space/cluster_c.data}
    \plotProbe[color=colorDirBA, style=dashed, width=thick]{-11}{-31.5}
    \plotDirection[color=colorDirBA]{-2.5}{-4.0}{-3.0}{1.5}
    \plotProbe[color=colorDirBC, style=dashed, width=thick]{0.5224}{-2.694}
    \plotDirection[color=colorDirBC]{-2.5}{-4.0}{4.2}{-0.5}
    \fill[colorB] (axis cs:-2.5,-4.0) circle (4pt);
  \end{axis}
  \node[anchor=north, yshift=-6pt] at (axisDirections.south)
    {\textbf{(b)} Interpretable directions};
  \end{scope}

  % ============================================================
  % (c) BOTTOM-LEFT — Sparse Autoencoder
  % ============================================================
  \begin{scope}[xshift=0pt, yshift=-\ovRowGap]
  \begin{axis}[
    name=axisSAE,
    width=\ovPanelWd, height=\ovPanelHt,
    xmin=-8.5, xmax=8.5, ymin=-6, ymax=9,
    axis lines=none, ticks=none, clip=true,
  ]
    \draw[colorA, thick, -{Stealth[length=5pt,width=4pt]}]
      (axis cs:0,0) -- (axis cs:0,7);
    \node[colorA, font=\small\bfseries, above] at (axis cs:0,7.2) {$A$};
    \draw[colorB, thick, -{Stealth[length=5pt,width=4pt]}]
      (axis cs:0,0) -- (axis cs:-6.062,-3.5);
    \node[colorB, font=\small\bfseries, below left] at (axis cs:-6.062,-3.7) {$B$};
    \draw[colorC, thick, -{Stealth[length=5pt,width=4pt]}]
      (axis cs:0,0) -- (axis cs:6.062,-3.5);
    \node[colorC, font=\small\bfseries, below right] at (axis cs:6.062,-3.7) {$C$};
    \plotPoints[color=colorA, mark=*]{figures/latent_space/sae_right_a.data}
    \plotPoints[color=colorB, mark=square*]{figures/latent_space/sae_right_b.data}
    \plotPoints[color=colorC, size=2pt, mark=triangle*]{figures/latent_space/sae_right_c.data}
  \end{axis}
  \node[anchor=north, yshift=-6pt] at (axisSAE.south)
    {\textbf{(c)} Sparse autoencoder};
  \end{scope}

% ============================================================
  % (d) BOTTOM-RIGHT — Clustering
  % ============================================================
  \begin{scope}[xshift=\ovPanelGap, yshift=-\ovRowGap]
  \begin{axis}[
    name=axisClustering,
    latent embed, width=\ovPanelWd, height=\ovPanelHt,
  ]
    \plotEllipse[color=colorA, width=very thick, fill=colorA, fillopacity=0.08]
      {-3}{1.5}{15}{1.4}{0.8}
    \plotEllipse[color=colorB, width=very thick, fill=colorB, fillopacity=0.08]
      {-2.5}{-4}{10}{1.3}{0.7}
    \plotEllipse[color=colorC, width=very thick, fill=colorC, fillopacity=0.08]
      {4.2}{-0.5}{-20}{0.9}{1.6}
    \plotPoints[color=colorA, mark=*]{figures/latent_space/cluster_a.data}
    \plotPoints[color=colorB, mark=square*]{figures/latent_space/cluster_b.data}
    \plotPoints[color=colorC, size=2pt, mark=triangle*]{figures/latent_space/cluster_c.data}
  \end{axis}
  \node[anchor=north, yshift=-6pt] at (axisClustering.south)
    {\textbf{(d)} Clustering};
  \end{scope}

\end{tikzpicture}
}

\caption{
  Overview of the four latent space analysis methods studied in this chapter, applied to the same point cloud.
}
\label{fig:latent_space_overview}
\end{figure}

The four families of methods reviewed in this chapter, namely probing, interpretable directions, sparse autoencoders, and clustering, are summarized in Figure~\ref{fig:latent_space_overview} and all build upon the \ConceptHypothesisName introduced in Section~\ref{sec:concept_hypothesis}.
Each of them approaches the question from a different angle, but they share a common objective: uncover how \gls{lr} are structured and what information can be extracted from them.

However, a fundamental difficulty remains.
There is currently no consensus on how to rigorously evaluate whether these methods truly capture the internal concepts of a model.
Most existing metrics assess reconstruction quality or linear separability, but these do not directly measure whether the discovered \gls{cs} correspond to meaningful abstractions from the model's perspective.
This leaves open the question of what it means, formally, for a concept to be represented in a \gls{nn}.

Addressing this question may require contributions from disciplines beyond computer science and mathematics.
As the notion of semantics is inherently tied to human interpretation, a purely formal or computational account may be insufficient.
Insights from philosophy, cognitive science, or psychology could prove valuable in grounding the notion of concept in a way that is both theoretically rigorous and empirically tractable.
Despite these open questions, the methods studied in this chapter provide a practical foundation for analyzing the internal representations of \gls{nn}.

Building on the \ConceptHypothesisName and the insights developed throughout this chapter, this manuscript investigates a representation-based approach to explainability in \gls{rl}.
More specifically, it builds on the intuition underlying deep clustering to identify meaningful structures within the representations learned by \gls{rl} agents.
To establish the necessary foundations, the next chapter introduces the principles of \gls{rl}.
The subsequent chapter then presents the proposed method for analyzing these representations and studying their potential for explainability in \gls{rl} (Chapter~\ref{chap:clustering_agent_representation}).